\emph{Druck:} Eine Person mit $m = \mSkO$ steht auf einer Kufe, die auf einer Länge von etwa $l = \lKO$ mit einer Breite von $b = \bKO$ auf dem Eis aufliegt:
\begin{align}
 p &= \pKF = \frac{\mSk\cdot\gE}{\lK\cdot\bK} = \pK \approx \result{\pKbP{0}{2}}
\end{align}

\emph{Schmelzpunkterniedrigung:} Weil Eis eine kleinere Dichte hat als Wasser, sinkt der Schmelzpunkt unter Druck. Die Gleichung von Clausius-Clapeyron gibt die Änderung pro Druckeinheit, mit $T_\text{f} = \TfO$, $L_\text{f} = \LfEO$, $\rho_\text{E} = \rhEO$ (FoTaBe S.~187, 190) und $\rho_\text{W} = \rhWnO$ bei $\vartheta_0 = \TnO$ (FoTaBe S.~192):
\begin{align}
 k = \frac{\Delta T}{\Delta p} &= \kCCF \\
   &= \frac{\Tf}{\LfE}\cdot\left(\frac{1}{\rhWn}-\frac{1}{\rhE}\right) = \kCC
\end{align}
Unter der Kufe sinkt der Schmelzpunkt deshalb nur um
\begin{align}
 \Delta T &= \dTKF = \kCC\cdot\pK = \dTK \approx \result{\dTKP{0}{2}}
\end{align}
Damit Eis von $\vartheta_\text{E} = \TEisO$ unter der Kufe schmilzt, müsste der Schmelzpunkt um $\vartheta_\text{E}-\vartheta_0 = \dTn$ sinken, also bräuchte es einen Druck von
\begin{align}
 p &= \pnF = \frac{\dTn}{\kCC} = \pn \approx \result{\pnbP{0}{2}}
\end{align}
also etwa \pnbP{0}{1} statt \pKbP{0}{1}. Die Druckschmelzung erklärt das Gleiten nicht. Die Kufe gleitet auf einer sehr dünnen, flüssigkeitsähnlichen Schicht an der Eisoberfläche, die durch Reibungswärme noch verstärkt wird.
